\capitulo{1}{Introducción}
En las últimas décadas, la incidencia del cáncer de piel, especialmente el melanoma, ha experimentado un crecimiento preocupante tanto en su aparición como en su mortalidad a nivel global. \cite{gutierrez2019trends} En España,  aunque la tasa de mortalidad de melanomas no se encuentra entre las más elevadas entre las diferentes causas, su evolución clínica, la agresividad clínica del tratamiento y su fuerte relación con factores ambientales aumentan la prioridad de su análisis ya sea en un contexto epidemiológico o en salud pública. Según la Asociación Española Contra el Cáncer\cite{CancerPielPresentaciona}en 2024 se estimaron 20.854 nuevos casos de  cánceres de piel de los cuales 6.070 de ellos son melanomas. 
  

Entre los principales factores de riesgo de desarrollo de un melanoma, caben destacar aquellos de carácter ambiental, siendo la radiacion ultravioleta uno de los ampliamente reconocidos social y cientificamente \cite{bozProspectiveCohortAnalysis2022a}. Sin embargo, tambien existen otras variables como la altitud \cite{delfioreAltitudeEffectCutaneous2022a}  y las temperaturas extremas \cite{millerHeat2024} que contribuyen al desarrollo y avance de esta enfermedad. 

En el contexto actual del cambio climático, el análisis geográfico y ambiental de la mortalidad por melanoma adquiere una especial relevancia, tanto para una mejor comprensión de los factores que influyen en su desarrollo como para el diseño de estrategias preventivas más eficaces. El progresivo aumento de las temperaturas asociado al calentamiento global, tal como advierte la Academia Española de Dermatología y Venereología \cite{CambioClimaticoCancer2015} , refuerza la necesidad de integrar variables climáticas en el estudio de esta patología.


Este Trabajo de Fin de Grado se centra en el desarrollo de una aplicación web interactiva y divulgativa orientada a la concienciación y prevención del cáncer de piel, en particular del melanoma maligno. La herramienta permite visualizar la variabilidad provincial de factores ambientales como la altitud, la radiación ultravioleta y  la temperatura máxima del día, integrando estas dos últimas variables en una base de datos propia que se actualiza diariamente, además de otra presentación de las defunciones por \textit{melanoma} provenientes del Instituto Nacional de Estadística. Además de ofrecer mapas interactivos y acceso a contenido educativo, la aplicación actúa como un sistema recomendador: alerta al usuario sobre condiciones de riesgo en su provincia y adapta los mensajes según su fototipo de piel. De forma complementaria, estima un posible riesgo acumulado de exposición para el día actual, considerando factores ambientales clave. En el contexto del cambio climático y el aumento de temperaturas, este tipo de herramienta cobra especial relevancia tanto por su valor divulgativo como por su potencial para generar una base de datos útil en futuras investigaciones. A largo plazo, se espera que la acumulación de información permita identificar patrones y asociaciones significativas entre condiciones climáticas y la incidencia del melanoma en la población.





